\subsection{TRANSPOSE OF A LINEAR MAPPING}

Let E, F be two left A-modules; for every linear mapping \(u: E \to F\) , the mapping \(\operatorname{Hom}(u, l_{A_s})\) is a linear mapping of the right A-module \(F^*\) into the right A-module \(E^*\) (§1, no. 2), called the transpose of u.

In other words:

DEFINITION 5. For every linear mapping \(u\) of an A-module \(E\) into an A-module \(F\) , the linear mapping \(y^{*} \mapsto y^{*} \circ u\) of the dual \(F^{*}\) of \(F\) into the dual \(E^{*}\) of \(E\) is called the transpose of \(u\) and is denoted by \(^{t}u\) .

The transpose \(^{t}u\) is therefore defined by the relation

\[
\langle u (x), y ^ {*} \rangle = \langle x, ^ {t} u (y ^ {*}) \rangle \quad \text { for   all } x \in \mathrm{E} \text { and   all } y ^ {*} \in \mathrm{F} ^ {*}.\tag{15}
\]

Definition 5 applies without alteration to right A-modules and is then equivalent to the relation

\[
\langle y ^ {*}, u (x) \rangle = \langle^ {t} u (y ^ {*}), x \rangle \quad \text { for   all } x \in \mathrm{E} \text { and   all } y ^ {*} \in \mathrm{F} ^ {*}.
\]

Formulae (9) and (10) of § 1, no. 2 here give

\[
{ } ^ { t } ( u _ { 1 } + u _ { 2 } ) = { } ^ { t } u _ { 1 } + { } ^ { t } u _ { 2 }\tag{16}
\]

for two elements \(u_{1}, u_{2}\) of \(\operatorname{Hom}_{\mathbf{A}}(\mathbf{E}, \mathbf{F})\) and

\[
{ } ^ { t } ( v \circ u ) = { } ^ { t } u \circ { } ^ { t } v\tag{17}
\]

for \(u \in \operatorname{Hom}_{\mathbf{A}}(\mathbf{E}, \mathbf{F})\) and \(v \in \operatorname{Hom}_{\mathbf{A}}(\mathbf{F}, \mathbf{G})\) , \(\mathbf{G}\) being a third A-module; finally, clearly

\[
{ } ^ { t } \mathbf { l } _ { \mathbf { E } } = \mathbf { l } _ { \mathbf { E } ^ { * } } .\tag{18}
\]

Remark. From (17) and (18) it follows that if \(u\) is left (resp. right) invertible, \(^{t}u\) is right (resp. left) invertible.

PROPOSITION 8. Let \(u: \mathbf{E} \to \mathbf{F}\) be an A-linear mapping, \(\mathbf{M}\) a submodule of \(\mathbf{E}\) and \(\mathbf{M}'\) the orthogonal of \(\mathbf{M}\) in \(\mathbf{E}^*\) ; the orthogonal of \(u(\mathbf{M})\) in \(\mathbf{F}^*\) is the inverse image \(^t u^{-1}(\mathbf{M}')\) .

This follows immediately from (15).

COROLLARY. The orthogonal of the image \(u(\mathbf{E})\) in \(\mathbf{F}^*\) is the kernel \({}^t u^{-1}(0)\) of \({}^t u\) .

The orthogonal of \(\mathbf{E}\) in \(\mathbf{E}^*\) is 0.

If \(u: \mathbf{E} \to \mathbf{F}\) is an isomorphism, \({}^{t}u: \mathbf{F}^{*} \to \mathbf{E}^{*}\) is an isomorphism and if \(v: \mathbf{F} \to \mathbf{E}\) is the inverse isomorphism of \(u\) , \({}^{t}v\) is the inverse isomorphism of \({}^{t}u\) (formulae (17) and (18)).

DEFINITION 6. Given an isomorphism u of an A-module E onto an A-module F, the transpose of the inverse isomorphism of u (equal to the inverse isomorphism of the transpose of u) is called the contragredient isomorphism of u and denoted by \(\check{u}\).

The isomorphism \(\check{u}\) is thus characterized by the relation

\[
\langle u (x), \check {u} (y ^ {*}) \rangle = \langle x, x ^ {*} \rangle \quad \text { for } x \in \mathbf {E}, x ^ {*} \in \mathbf {E} ^ {*}.\tag{19}
\]

If \(v: \mathbf{F} \to \mathbf{G}\) is an isomorphism, the contragredient isomorphism of \(v \circ u\) is \(v' \circ \check{u}\) .

In particular, the mapping \(u \mapsto \check{u}\) is an isomorphism of the linear group \(\mathbf{GL}(\mathbf{E})\) onto a subgroup of the linear group \(\mathbf{GL}(\mathbf{E}^*)\) .

Let \(\sigma: \mathbf{A} \to \mathbf{B}\) be an isomorphism of a ring \(\mathbf{A}\) onto a ring \(\mathbf{B}\) , \(\mathbf{E}\) a left \(\mathbf{A}\) -module, \(\mathbf{F}\) a left \(\mathbf{B}\) -module and \(u: \mathbf{E} \to \mathbf{F}\) a semi-linear mapping (§ 1, no. 13) relative to \(\sigma\) . Let \(\sigma^{-1}\) be the inverse isomorphism of \(\sigma\) ; for all \(y^{*} \in \mathbf{F}^{*}\) , the mapping \(x \mapsto \langle u(x), y^{*} \rangle^{\sigma - 1}\) of \(\mathbf{E}\) into \(\mathbf{A}\) is a linear form; if it is also denoted by \(^t u(y^*)\) , a mapping \(^t u: \mathbf{F}^{*} \to \mathbf{E}^{*}\) is defined which is also called the transpose of the semi-linear mapping \(u\) ; it is thus characterized by the identity

\[
\langle u (x), y ^ {*} \rangle = \langle x, ^ {t} u (y ^ {*}) \rangle^ {\sigma}\tag{20}
\]

for \(x \in \mathbf{E}\) , \(y^* \in \mathbf{F}^*\) . It is immediately verified that \({}^t u\) is a semi-linear mapping relative to \(\sigma^{-1}\) . If \(v\) denotes the mapping \(u\) considered as an A-linear mapping of \(\mathbf{E}\) into \(\sigma_*(\mathbf{F})\) (§ 1, no. 13), we may write \(u = \phi \circ v\) , where \(\phi\) is the identity mapping \(\sigma_*(\mathbf{F}) \to \mathbf{F}\) , considered as a semi-linear mapping relative to \(\sigma\) . It is immediate that \({}^t u = {}^t v \circ {}^t \phi\) and \(({}^t \phi, \sigma^{-1})\) is a di-isomorphism of \(\mathbf{F}^*\) onto \((\sigma_*(\mathbf{F}))^*\) relative to the isomorphism \(\sigma^{-1}\) ; this relation allows us immediately to extend the properties of transposes of linear mappings to transposes of semi-linear mappings.

